% ============================================================================
\section{When Can a Shared Update Make Joint Progress?}
\label{sec:theory}
% ============================================================================

\subsection{Expected progress under the chosen task mixture}

Condition on the training state, freeze the rollout distributions in
$L_i^t$, and choose counts before sampling. Let
$\mu_i=\nabla L_i^t(\theta_t)$ and
$\Sigma_i=\operatorname{Cov}(g_{i,s})$. For the following local analysis,
micro-batches are independent and have finite second moments. The
diagnostic experiment uses independent draws; the training streams require
the finite-population qualification in Appendix~\ref{app:local_objective}.

Let $D$ be a positive diagonal scaling fixed before the step. For SGD,
$D=I$; for Adam, we use
$D=\operatorname{diag}((\sqrt{\widehat v_{\rm pre}}+\epsilon)^{-1})$,
where $\widehat v_{\rm pre}$ is the pre-step bias-corrected second moment.
For the idealized update $\theta'=\theta_t-\eta DA$, define
\begin{equation}
    \mu_w=\sum_jw_j\mu_j,\qquad
    K_{ij}=\mu_i^\top D\mu_j,\qquad
    a_i=(Kw)_i=\mu_i^\top D\mu_w.
    \label{eq:transfer_margins}
\end{equation}
The expected first-order change in task $i$'s loss is $-\eta a_i$.
Thus all $a_i>0$ means that the chosen mean update is a common descent
direction. A small positive $a_i$ leaves little margin against sampling
noise. If $a_i\leq0$, changing counts at fixed weights cannot produce
strictly positive expected first-order progress for that task.

This last statement concerns this mixture at this checkpoint. It does
not rule out another mixture, another update direction, or later recovery.
Pairwise negative gradient cosines do not determine the signs of $Kw$.
The relevant interaction uses one $D$; cosines between $D\mu_i$ and
$D\mu_j$ instead use $D^2$. Actual AdamW includes momentum, an updated
second moment, clipping, and weight decay. Its branch updates must be
measured separately from this frozen-scaling approximation.

\subsection{Sampling uncertainty relative to task progress}

Let $X_i=\mu_i^\top DA$ be the sampled first-order progress. Its variance is
\begin{equation}
    v_i(m)=\operatorname{Var}(X_i)
      =\sum_j\frac{w_j^2}{m_j}
          \mu_i^\top D\Sigma_jD\mu_i.
    \label{eq:directional_variance}
\end{equation}
For $a_i>0$, the one-sided Chebyshev--Cantelli inequality gives
\begin{equation}
    \Pr(X_i\leq0)\leq\frac{v_i(m)}{v_i(m)+a_i^2},\qquad
    \Pr(\exists i:X_i\leq0)
      \leq\min\!\left\{1,\sum_i\frac{v_i(m)}{v_i(m)+a_i^2}\right\},
    \label{eq:directional_reliability}
\end{equation}
where the joint bound requires every $a_i>0$. Reliability therefore
depends on noise relative to the progress margin, not on raw noise alone.
These bounds require no Gaussian assumption, but can be loose. Replacing
population quantities with noisy estimates does not produce a statistical
certificate. Appendix~\ref{app:joint_progress} proves the bounds;
Appendices~\ref{app:budget_reliability} and~\ref{app:descent} give a
sufficient sampling budget and the additional curvature conditions for
finite-step descent. Directional sampling tests and stochastic
multiobjective sampling have established precedents
\citep{bollapragada2018adaptive,zhao2024adaptive}; our empirical question
is whether these quantities predict MOPD regressions and useful interventions.

\subsection{GPAS controls a conservative variance bound}

GPAS measures total gradient variation after optimizer scaling:
\begin{equation}
    e_i=\mathbb E\|D(g_i-\mu_i)\|_2^2,\qquad
    V(m)=\mathbb E\|D(A-\mu_w)\|_2^2
         =\sum_i\frac{w_i^2e_i}{m_i}.
    \label{eq:batch_variance}
\end{equation}
It controls each directional variance through
\begin{equation}
    v_i(m)\leq\|\mu_i\|_2^2 V(m).
    \label{eq:trace_directional_bound}
\end{equation}
This is a conservative surrogate: fluctuations orthogonal to a task's
gradient contribute to $V$ without affecting its first-order progress.
Appendix~\ref{app:trace_counterexample} gives aligned task means for
which trace-optimal allocation has worse directional reliability than
another allocation. GPAS is therefore an intervention to test, not an
optimal rule for joint-progress probability.

Ignoring integer bounds, minimizing $V$ at $\sum_i m_i=G$ yields
\begin{equation}
    m_i\propto w_i\sqrt{e_i}.
    \label{eq:batch_gpas_allocation}
\end{equation}
This is classical Neyman allocation
\citep{neyman1934representative,cochran1977sampling} in the optimizer's
coordinates. The four-task implementation enumerates all feasible
integer counts and minimizes the estimated $V$ exactly over that set;
Appendix~\ref{app:allocation_proofs} gives the derivation. Under the
study's bounds, $m_i\leq8$ versus Uniform's $m_i=4$ implies
$V_{\rm Uniform}/V(m)\leq2$ when the denominator is positive.
This limits variance reduction under the local model, not training speed.

\subsection{Online estimates and measured computation}
\label{sec:online_estimates}

GPAS reuses the unweighted micro-batch gradients. With task mean $\bar g_i$,
\begin{equation}
    \widehat e_i=\frac1{m_i-1}\sum_{s=1}^{m_i}
                    \|D(g_{i,s}-\bar g_i)\|_2^2.
    \label{eq:noise_estimate}
\end{equation}
The next step uses the moving average
$\widetilde e_i\leftarrow\rho\widetilde e_i+(1-\rho)\widehat e_i$.
A running mean and scalar squared-deviation sum avoid retaining all
micro-batch gradients; one parameter-sized mean buffer is reused across
task blocks. The first step uses Uniform. We measure the additional
reductions, memory, and time rather than assuming negligible overhead.

For marginal micro-batch time $\tau_i$ and fixed step time $C$, the
cost-aware mode minimizes
\begin{equation}
    J_C(m)=T(m)V(m)
      =\left(C+\sum_i m_i\tau_i\right)
         \sum_i\frac{w_i^2e_i}{m_i}.
    \label{eq:batch_cost_objective}
\end{equation}
This criterion measures time times variance; it does not generally equal
learning progress per second. The continuous rule without count bounds at
$C=0$ is $m_i\propto w_i\sqrt{e_i/\tau_i}$, and the finite implementation
enumerates feasible counts. The timing model is fitted to measured task
times and checked against end-to-end GPU hours. Directional allocation
is a conditional extension in Appendix~\ref{app:directional_extension},
to be investigated only if trace noise fails to predict the uncertainty
that affects task progress.

\begin{algorithm}[t]
\caption{One training step with GPAS.}
\label{alg:batch_gpas}
\small
\begin{minipage}{0.97\linewidth}
\begin{enumerate}
    \item Choose feasible integer counts from the previous noise and time
    averages; use Uniform on the first step. Read the pre-step scaling $D$.
    \item For each task, generate $m_i$ micro-batches and obtain its
    teacher's top-$k$ scores on the student responses.
    \item Compute the dense loss, accumulate unweighted task gradients,
    record their variation under $D$, and add the task mean to $A$ with
    weight $w_i$.
    \item Update the noise and time averages, then apply the usual
    optimizer step with gradient $A$.
\end{enumerate}
\end{minipage}
\end{algorithm}
